\textbf{Status of this document:} Independent specification draft, written in
Internet-Draft style, not submitted to the IETF. It has no standing in any
standards process, and the RFC~2119 and RFC~8174 key words are used as in an
Internet-Draft.

\medskip
\noindent\textbf{How to read this document.}
The comparisons in Sections~\ref{sec:related-work} and~\ref{sec:mapping}
are the author's reading of the cited drafts at the versions listed in
\texttt{paper/sources/MANIFEST.md}. Where this document and a cited
draft differ, the cited draft is authoritative for its own scope; this
document is authoritative only for the Bridge Contract. The phrase
``as read'' marks an interpretation of a cited draft, not a fact about
it. The ``Why this obligation is out of <draft>'s scope'' paragraphs in
Section~\ref{sec:related-work} argue that, in the versions read, the gap is
one of scope rather than an oversight. The wire format (Appendix~\ref{app:wire-format}) is
normative for client--bridge messages; the unit of conformance is the
committed event log and the final session state. The limits of the
evidence are stated in Section~\ref{sec:limitations} and the
limitations register (Table~\ref{tab:limitations}).

% ─────────────────────────────────────────────────────────────
% Abstract and keywords. The block below is regenerated from
% metadata.yaml by `make metadata`; edit metadata.yaml, not this file.
% ─────────────────────────────────────────────────────────────
% ABSTRACT:START
\begin{abstract}
\noindent
Existing agent protocols --- the Agent Transfer Protocol (AGTP), the
Governance Audit Record (GAR) and the Protocol for Agent Communication
Exchange (PACE) --- define transport, audit and resumption semantics
for agent sessions, but, in the versions read, none defines what an intermediate bridge MUST
do when the client is detached and the agent keeps executing, or when
the human in the loop is unavailable. We argue that the reason is
scope rather than oversight; this rests on a reading of those drafts
and would weaken if one of them were extended to own these duties.

This document defines the Bridge Contract BC = (S, E, R, C): a Session
Identity Descriptor, Mandatory Governance Events, a Resume Record that
the bridge keeps for itself, and Conformance Criteria. Re-attach is a
challenge-response in which the client proves possession of a
registered key over a single-use nonce and the last event it saw,
and the bridge accepts only a consistent prefix. The log is
hash-chained, which detects edits and reordering inside it but not
the truncation of its last events without a head held outside the log. The normative model
is a product state machine of attachment, human and lifecycle, with
invariants I1--I15 and derived timers. The contract complements
AGTP, GAR and PACE and does not replace them.

Evidence: the product model is checked with TLC on bounded
configurations; a conformance suite (T1--T158) and a mutation check
in which every MUST invariant except the inspection-only I-SIGN-JWKS has
a killed mutant, on a Python bridge; vectors V01--V134 pass on it, and
all but the checkpoint-emission vectors (V52--V54c, V88, V89; they need a
feature the crate lacks) pass on a Rust crate. Both implementations have the
same author; no third-party implementation exists. The evidence
shows the consistency of one reading of the text, not independent
confirmation; the limits are stated in the limitations section and
the limitations register.

\end{abstract}

\bigskip
\noindent\textbf{A note on evidence.} This document arose from the development of Termenbox,
the author's proprietary, free-to-use project, a self-hosted web terminal for AI coding agents, which may adopt
the specification. The specification is implementation-agnostic.


\medskip
\noindent\textbf{Keywords:} agent, bridge contract, session governance, AGTP, GAR, PACE, conformance, session resumption, specification draft.
% ABSTRACT:END

\section*{Normative summary}
\addcontentsline{toc}{section}{Normative summary}

This page restates the obligations of a conformant bridge; the sections
cited are authoritative where the wording differs.

\smallskip\noindent\textbf{State.} A live session is in a cell
$(a,h)$, $a \in \{\textsc{Att}, \textsc{Det}, \textsc{Res}\}$ (client
attached, detached, re-attaching), $h \in \{\textsc{None},
\textsc{Pend}\}$ (no escalation, or the agent waits on a human), or in the
terminal cell $\bot$. The transition table
(Section~\ref{sec:product-model}) is normative; the names ATTACHED,
DETACHED, PENDING\_HUMAN, RESUMING and TERMINATED are a derived view. Every
transition emits a governance event or is an explicit idempotent no-op.

\smallskip\noindent\textbf{A bridge MUST} (Sections~\ref{sec:spec}
and~\ref{sec:state-machine}):
\begin{itemize}[leftmargin=1.5em,itemsep=0pt,topsep=2pt]
\item append every event, hash-chained (I5), durably and \emph{before}
the cell changes or the client is acknowledged (I7, I12), and rebuild each
session after a restart by replaying its log, as detached
(Section~\ref{sec:recovery});
\item emit \texttt{DETACHMENT} and \texttt{RESUME\_PREPARED} before
entering a detached cell (I1), and queue, never forward, tool calls while
$h=\textsc{Pend}$ (I2); queued calls carry idempotency keys, the queue
is bounded, calls and human approvals expire, and drain revalidates the
current scope (I13); a call is forwarded only if an approving answer
whose reported call hash (tool, permission and content) equals the hash recorded at queue time
exists (I14);
\item let authority only narrow, with a version counter (I10); refuse a
widening it is asked to apply, and report a widening the runtime
announces as \texttt{AUTHORITY\_EXTENSION\_REJECTED} with source
\texttt{observed} (I3; the bridge prevents the first and can only report the
second);
\item treat attachment as derived from the set of connected client keys
(I15): a further client joining or leaving while another is connected
changes no cell;
\item accept a re-attach only as a challenge-response: a signature under a
registered client key over a fresh, single-use nonce
and the last event the client saw, where that event is in the bridge's
chain (I4, I9, I11); any other attempt is rejected without changing state;
emit \texttt{RESUME\_VERIFIED} on $\textsc{Res}\to\textsc{Att}$ after
re-verifying the chain;
\item derive the timers from one value, $T_{\text{detach}}$, with
$T_{\text{call}} \le T_{\text{human}}$ (MUST) and
$T_{\text{human}} \le T_{\text{detach}}$ as the default and a recommendation
(default 30~min, profile \emph{persistent} 7~days),
time out a pending human and a detachment, and keep $\bot$ terminal (I6).
\end{itemize}
It SHOULD sign events in GEC-less deployments (I-SIGN-EVENT) and emit
signed checkpoints anchored in an external transparency log
(I-CHECKPOINT). It MUST NOT accept a JWKS endpoint from session metadata
(I-SIGN-JWKS).

\smallskip\noindent\textbf{Assumptions and scope.} The bridge is
trusted (A4), the clients' private keys are not stolen (A7), and the
durable log honours append (A1). A Byzantine client without a registered
key and a restart of a single bridge are in scope; replication, partial
detach, a Byzantine bridge and multi-hop composition are not
(Section~\ref{sec:scope}). I2 and I13 hold for calls the runtime submits
to the bridge (Section~\ref{sec:coverage}).

\smallskip\noindent\textbf{Evidence} (Section~\ref{sec:conformance},
Appendix~\ref{app:formal}). TLC finds no violation of the product model on
bounded configurations (the two-session run never fires Disconnect,
Resolve, Submit or Drain; the detach paths are checked only in
single-session runs) and kills every mutant of the model; the pytest
cases and the mutants (every invariant I1--I15 has a killed mutant) cover
the Python reference bridge; JSON vectors V01--V134 run on it, and all but those that expect checkpoint events run on a Rust crate. No third-party implementation exists yet.

\tableofcontents
\newpage
